% conclusion.tex

\section{Discussion}
\label{sec:discussion}

\paragraph{Topic shortcut and label equivalence.}
Three observations argue against a topic shortcut: unstable topic OOD (F1\,=\,0.541$\pm$0.207), scrambled labels within 3.9\,pp, and label-free MLM detection (F1\,$\geq$\,0.97).

\paragraph{Single-LLM confound.}
ToxicChat (+18--42\,pp; Table~\ref{tab:human_data}) and cross-LLM experiments (Appendix~\ref{app:cross_llm}) confirm transfer. Mixed-LLM 9-class maintains 0.818 while MLM collapses; MLM may capture LLM-specific surface patterns while classification extracts style-invariant features.

\paragraph{Variance reversal on human data.}
Classification supervision is associated with lower variance on synthetic OOD (4--20$\times$ observed range, 3 seeds); on ToxicChat this reverses (std\,=\,0.087 vs.\ 0.013). Frozen probing localizes MLM instability to Stage-2 (Appendix~\ref{app:frozen_verify}).

\paragraph{Encoder capacity and lexical analysis.}
Three encoders (125M--355M) differ in pretraining; same-architecture scaling is needed. TF-IDF collapses on ActorAttack (0.261 vs.\ 0.961 on DD).

\paragraph{Regularization mechanism.}
Frozen probing near-parity (AA gap 3.1\,pp), label irrelevance ($d$\,=\,0.09), and mean pooling matching BiGRU ($|\Delta| < 1$\,pp) converge: classification DAPT provides robustness through optimization regularization~\citep{chen2020simclr,grill2020byol}, not semantic feature learning.

\section{Conclusion}

Classification DAPT provides robustness through optimization regularization, not semantic learning, and is 266$\times$ faster than LLM-as-judge.

\section*{Limitations}
\label{sec:limitations}

\paragraph{Multi-turn human validation.}
All multi-turn findings are derived from synthetic conversations generated by Qwen3-8B. ToxicChat validates transfer to human-authored single-turn attacks (+18--42\,pp over vanilla), but no human-generated multi-turn jailbreak dataset exists for evaluation. Whether the evidence hierarchy and classification supervision advantages generalize to human-authored multi-turn attacks remains the primary open question.

\paragraph{Turn-count distribution.}
Benign conversations contain 12--20 turns, DD jailbreaks 6--10, and ActorAttack jailbreaks 3--8; the three categories have non-overlapping turn counts. Three observations argue against length as the primary detection cue: (a) early detection at equal-length prefixes $K$=1 and $K$=2 achieves AA OOD F1\,=\,0.773 and 0.935 (Appendix~\ref{tab:actorattack_full}); (b) vanilla models access the same length signal but achieve near-zero OOD recall; and (c) truncating benign conversations to 3 turns yields FPR\,=\,0\% (Appendix~\ref{app:length_independence}).

\paragraph{Temporal aggregation dilution.}
When subtle attacks (ActorAttack) are followed by extended benign padding, BiGRU temporal aggregation dilutes attack signals (F1 drops from 0.963 to 0.383; Appendix~\ref{app:length_independence}); in practice, adversaries would not dilute their own successful attacks.

\paragraph{OOD benign distribution.}
OOD results reflect attack-side generalization only (Section~\ref{sec:discussion}). Even with matched benign distributions, residual calibration error (ECE\,=\,0.126--0.141; Appendix~\ref{app:calibration}) suggests temperature scaling is needed for deployment.

\paragraph{Frozen encoder assumption.}
The frozen Stage-1 design isolates representation quality but may underestimate full-pipeline performance. End-to-end fine-tuning closes this gap to $\leq$5.8\,pp on ActorAttack (F1\,=\,0.997$\pm$0.005 vs.\ 0.939$\pm$0.047 frozen; Table~\ref{tab:e2e}; Appendix~\ref{app:frozen_verify}).

\paragraph{Dataset scale and cross-LLM transfer.}
Scale validation (1{,}000 conversations; Appendix~\ref{app:scale_validation}) confirms the hierarchy holds (DAPT advantage 40--56\,pp over vanilla). Cross-LLM validation shows jailbreak detection transfers when benign distribution is held constant (F1 recovers to 0.800--1.000), but benign-side F1 drops to 0.370--0.430 with unfamiliar benign text (Appendix~\ref{app:cross_llm}). The perturbation evaluation ($n$=76) remains underpowered.

\paragraph{Strategy labels.}
Persuasion strategy labels ($n$=50 human validation, $\kappa$=0.78) are produced by Qwen3-8B; the validation sample yields small per-class sizes across 9 categories, limiting reliability estimates for individual strategies. OOD strategy accuracy degrades to 30.7\% ($\kappa$=0.126); interpretability is limited to known attack families. We plan to release annotation prompts and data for independent verification.

\paragraph{Detection architecture.}
Binary vs.\ 9-class equivalence (F1 gap $\leq$3.9\,pp) suggests the binary head is not the critical differentiator. FITD provides minimal discriminative signal (all variants F1\,$\geq$\,0.990), serving as confirmation that simple attacks do not challenge any variant.

\paragraph{Encoder scale.}
Whether findings hold at LLM-scale encoders (7B+) remains open; same-architecture scaling (e.g., DeBERTa-base vs.\ DeBERTa-large) is needed to attribute the capacity effect to scale alone.

\paragraph{Scope and single-turn transfer.}
Only persuasion-mediated multi-turn attacks fall within scope; single-turn token-level attacks~\citep{zou2023gcg}, non-persuasive strategies, and non-text modalities are not addressed. ToxicChat performance reflects zero-shot single-turn transfer; doubling the training corpus drops ToxicChat macro F1 from 0.596 to 0.341$\pm$0.004 while multi-turn OOD remains stable, confirming the architecture targets multi-turn detection.

\paragraph{Computation cost.}
Total pipeline cost for all seven variants is under 4 GPU-hours on a single A100.

\section*{Ethics Statement}

This work is defensive: it improves detection of multi-turn jailbreak attacks rather than developing new attacks. We acknowledge dual-use risk, as the persuasion taxonomy could inform adversarial actors. We mitigate this by using only synthetic conversations (no real harmful content), withholding attack-generation prompts, and focusing contributions on detection mechanisms. No human subjects were involved; all data will be released under CC BY-NC 4.0.
